\documentclass[11pt]{article}

\usepackage[margin=1in]{geometry}
\usepackage{booktabs}
\usepackage{amsmath}
\usepackage{amssymb}
\usepackage{array}
\usepackage{longtable}
\usepackage[hidelinks]{hyperref}
\usepackage{titlesec}
\titlespacing*{\section}{0pt}{1.4ex plus .2ex}{0.7ex}
\titlespacing*{\subsection}{0pt}{1.1ex plus .2ex}{0.4ex}
\titlespacing*{\paragraph}{0pt}{0.9ex plus .2ex}{1em}

\newcommand{\code}[1]{\texttt{#1}}
\newcolumntype{L}[1]{>{\raggedright\arraybackslash}p{#1}}
\newcolumntype{R}[1]{>{\raggedleft\arraybackslash}p{#1}}
\setlength{\emergencystretch}{3em}

\title{\textbf{ANCOMBC 2.15.1: computational changes and verification}}
\author{Package maintainer notes}
\date{September 17, 2026}

\begin{document}
\maketitle

\section{Scope}

Version 2.15.1 reduces the running time and the peak memory of the five
user-facing entry points \code{ancom}, \code{ancombc}, \code{ancombc2},
\code{secom\_linear} and \code{secom\_dist}. It adds the argument
\code{conservative} to \code{ancombc2}, which selects the rule used in the
sensitivity analysis to pseudo-count addition. It corrects two defects: the
failure of \code{ancombc} with \code{pseudo = 0} on a table containing zero
counts, and the failure of \code{ancombc2} with \code{pseudo\_sens = TRUE},
\code{trend = TRUE} and \code{global = FALSE}.

Three constraints apply to the computational changes. The implementation is R
only, with no compiled code. No package is added to \code{Imports},
\code{Depends} or \code{Suggests}. The numerical results of version 2.13.2 are
reproduced on every workload of the verification harness, at the tolerance
stated in Section~\ref{sec:verification}.

\section{Verification procedure}
\label{sec:verification}

The harness is \code{benchmarks/harness}, with five scripts:
\code{workloads.R} defines the workloads, \code{run.R} times them and records
peak memory, \code{compare.R} compares two sets of results,
\code{edge\_cases.R} exercises configurations that the workloads do not reach,
and \code{repro\_pseudo0.R} reproduces the \code{pseudo = 0} defect.

The harness contains 17 workloads. Twelve are the calls of the vignettes
\code{ANCOM.Rmd}, \code{ANCOMBC.Rmd}, \code{ANCOMBC2.Rmd} and \code{SECOM.Rmd},
reproduced with their arguments, their seed \code{set.seed(123)} and their
\code{n\_cl} values. Five are larger configurations at the Genus or OTU level.
Table~\ref{tab:workloads} lists them. Sizes are the taxa and samples that enter
the analysis after \code{prv\_cut} and \code{lib\_cut}.

All measurements were taken on one machine, an Apple M2 Max running R 4.5.3,
with the two installations exercised back to back. Vignette-scale times are
medians of three runs; large-scale times are single runs.

\begin{table}[htbp]
\centering
\scriptsize
\setlength{\tabcolsep}{4pt}
\caption{The 17 harness workloads. The last column is the 2.13.2 elapsed time.}
\label{tab:workloads}
\begin{tabular}{L{4.4cm} L{1.8cm} L{1.0cm} R{1.4cm} L{4.4cm} R{0.9cm}}
\toprule
Workload & Data & Level & Taxa $\times$ samples & Arguments & 2.13.2 (s) \\
\midrule
\multicolumn{6}{l}{\emph{Vignette scale}} \\
\code{ancom\_atlas} & atlas1006 & Family & $21 \times 873$ &
  \code{main\_var = "bmi"}, \code{adj\_formula = "age + region"},
  \code{n\_cl = 2} & 2.84 \\
\code{ancom\_dietswap} & dietswap & Family & $19 \times 222$ &
  \code{main\_var = "group"}, \code{rand\_formula = "(timepoint | subject)"},
  \code{n\_cl = 2} & 6.14 \\
\code{ancombc\_atlas} & atlas1006 & Family & $21 \times 873$ &
  \code{formula = "age + region + bmi"}, \code{global = TRUE},
  \code{n\_cl = 1} & 1.89 \\
\code{ancombc2\_qmp} & QMP simulation, $n = 150$ & OTU & $91 \times 143$ &
  \code{fix\_formula = "cont\_cov + cat\_cov"}, \code{pairwise = TRUE},
  \code{pseudo\_sens = TRUE}, \code{n\_cl = 2} & 7.13 \\
\code{ancombc2\_atlas\_perm} & atlas1006, permuted \code{bmi} & Genus &
  $118 \times 873$ & \code{fix\_formula = "bmi"}, \code{prv\_cut = 0},
  remaining arguments at their defaults & 10.63 \\
\code{ancombc2\_atlas\_family} & atlas1006 & Family & $21 \times 873$ &
  \code{fix\_formula = "age + region + bmi"}, global, pairwise, Dunnett and
  trend tests, \code{n\_cl = 2} & 26.84 \\
\code{ancombc2\_atlas\_interact} & atlas1006 & Family & $21 \times 873$ &
  \code{fix\_formula = "age * bmi + region"}, \code{n\_cl = 2} & 24.99 \\
\code{ancombc2\_dietswap} & dietswap & Family & $19 \times 222$ &
  \code{rand\_formula = "(timepoint | subject)"}, global, pairwise, Dunnett and
  trend tests, \code{n\_cl = 2} & 231.93 \\
\code{secom\_linear\_atlas} & atlas1006 & Phylum & $6 \times 873$ &
  \code{method = "pearson"}, \code{thresh\_len = 20}, \code{n\_cv = 10},
  \code{n\_cl = 2} & 1.38 \\
\code{secom\_dist\_atlas} & atlas1006 & Phylum & $6 \times 873$ &
  \code{R = 1000}, \code{n\_cl = 2} & 7.49 \\
\code{secom\_linear\_2groups} & atlas1006, regions CE and NE & Phylum &
  $6 \times 578$, $6 \times 181$ & \code{thresh\_len = 20}, \code{n\_cv = 10},
  \code{n\_cl = 2} & 1.42 \\
\code{secom\_dist\_2groups} & atlas1006, regions CE and NE & Phylum &
  $6 \times 578$, $6 \times 181$ & \code{R = 1000}, \code{n\_cl = 2} & 4.84 \\
\midrule
\multicolumn{6}{l}{\emph{Large scale}} \\
\code{ancom\_atlas\_genus\_large} & atlas1006 & Genus & $114 \times 873$ &
  \code{main\_var = "bmi"}, \code{adj\_formula = "age + region"},
  \code{n\_cl = 2} & 10.70 \\
\code{ancombc\_atlas\_genus\_large} & atlas1006 & Genus & $114 \times 873$ &
  \code{formula = "age + region + bmi"}, \code{global = TRUE},
  \code{n\_cl = 1} & 10.37 \\
\code{ancombc2\_atlas\_genus\_large} & atlas1006 & Genus & $114 \times 873$ &
  global, pairwise, Dunnett and trend tests, \code{B = 10} for the trend
  bootstrap, \code{n\_cl = 2} & 29.57 \\
\code{ancombc2\_qmp\_otu\_large} & QMP simulation, $n = 900$ & OTU &
  $91 \times 824$ & global, pairwise and Dunnett tests,
  \code{struc\_zero = TRUE}, \code{n\_cl = 2} & 9.04 \\
\code{secom\_linear\_atlas\_genus\_large} & atlas1006 & Genus &
  $114 \times 873$ & \code{prv\_cut = 0.10}, \code{thresh\_len = 100},
  \code{n\_cv = 10}, \code{n\_cl = 2} & 7.95 \\
\bottomrule
\end{tabular}
\end{table}

\paragraph{Equality criterion.} Let $a$ and $b$ denote the result objects of the
two versions for one workload. Components that carry an evaluation environment
are replaced by a sentinel before comparison. The workload is recorded as equal
when
\begin{equation}
\texttt{all.equal}(a, b, \texttt{tolerance} = 10^{-8}) = \texttt{TRUE},
\end{equation}
that is, when every pair of corresponding numeric components satisfies
$\overline{|a - b|} \, / \, \overline{|a|} \le 10^{-8}$ and every pair of
non-numeric components is identical. The comparison is repeated at tolerance
$0$, which reports bitwise equality.

\paragraph{Determinism.} Repeated runs of 2.13.2 on the 12 vignette workloads
are bitwise identical, including the workloads with \code{n\_cl = 2}, whose
parallel streams are generated by \code{doRNG}.

\paragraph{Edge cases.} \code{edge\_cases.R} covers nine configurations: a
non-finite response, a taxon whose usable samples cover one level of the group
factor, an all-zero taxon, a data-dependent model term, an unused level of the
group factor with and without \code{neg\_lb}, a missing group label, a
two-level group, and a covariate with missing values. Thirteen further
\code{ancom} configurations vary the main variable type, the adjustment
formula, the random-effects formula and \code{n\_cl}. All are bitwise identical
between the two versions, with the single exception of
\code{ancombc(pseudo = 0)}, which is the defect described in
Section~\ref{sec:bugs}.

\paragraph{Tests.} The package tests are the \code{testthat} files for
\code{ancom}, \code{ancombc}, \code{ancombc2} and \code{secom}, which include a
test of \code{conservative = FALSE}. They are run by \code{R CMD check} on the
built tarball, together with the five vignettes.

\paragraph{Memory measurement.} Peak memory is the increase in the \code{gc()}
``max used'' figure of the master process, summed over \code{Ncells} and
\code{Vcells}, between \code{gc(reset = TRUE)} before the call and \code{gc()}
after it. Workloads with \code{n\_cl > 1} spawn PSOCK worker processes whose
allocations \code{gc()} does not observe. For those workloads the reported
figure is a lower bound on the memory footprint of the run.

\section{Results}

Table~\ref{tab:timings} reports the back-to-back comparison of the 17
workloads. All 17 are equal at tolerance $10^{-8}$. Fourteen are bitwise
identical.

\begin{table}[htbp]
\centering
\footnotesize
\setlength{\tabcolsep}{5pt}
\caption{Elapsed time, peak master-process memory and equality for the 17
workloads.}
\label{tab:timings}
\begin{tabular}{l r r r r r l}
\toprule
Workload & 2.13.2 (s) & 2.15.1 (s) & Speedup & 2.13.2 (Mb) & 2.15.1 (Mb) &
Equality \\
\midrule
\code{ancom\_atlas}                     &   2.84 &  1.37 &  2.08 &  58.2 &  45.3 & bitwise \\
\code{ancom\_atlas\_genus\_large}       &  10.70 &  1.42 &  7.54 &  79.3 &  33.3 & bitwise \\
\code{ancom\_dietswap}                  &   6.14 &  3.48 &  1.76 &  48.8 &  50.3 & bitwise \\
\code{ancombc\_atlas}                   &   1.89 &  0.49 &  3.85 & 345.4 & 342.3 & bitwise \\
\code{ancombc\_atlas\_genus\_large}     &  10.37 &  0.76 & 13.68 & 363.9 & 341.5 & bitwise \\
\code{ancombc2\_atlas\_family}          &  26.84 &  3.41 &  7.87 & 294.1 & 111.2 & bitwise \\
\code{ancombc2\_atlas\_genus\_large}    &  29.57 &  3.61 &  8.19 & 294.0 & 168.4 & bitwise \\
\code{ancombc2\_atlas\_interact}        &  24.99 &  3.49 &  7.17 & 219.6 & 111.0 & bitwise \\
\code{ancombc2\_atlas\_perm}            &  10.63 &  0.89 & 12.01 & 369.3 & 211.9 & bitwise \\
\code{ancombc2\_dietswap}               & 231.93 & 64.63 &  3.59 & 356.3 & 309.5 & bitwise \\
\code{ancombc2\_qmp}                    &   7.13 &  1.85 &  3.85 & 342.0 & 133.0 & bitwise \\
\code{ancombc2\_qmp\_otu\_large}        &   9.04 &  2.13 &  4.24 & 362.7 & 152.0 & bitwise \\
\code{secom\_dist\_2groups}             &   4.84 &  3.52 &  1.38 & 128.2 & 116.5 & bitwise \\
\code{secom\_dist\_atlas}               &   7.49 &  6.06 &  1.24 &  79.2 &  69.6 & bitwise \\
\code{secom\_linear\_2groups}           &   1.42 &  0.35 &  4.09 & 125.6 & 119.1 & $10^{-8}$ \\
\code{secom\_linear\_atlas}             &   1.38 &  0.30 &  4.67 &  77.9 &  73.7 & $10^{-8}$ \\
\code{secom\_linear\_atlas\_genus\_large} & 7.95 &  0.44 & 17.90 & 118.6 &  91.2 & $10^{-8}$ \\
\bottomrule
\end{tabular}
\end{table}

The three \code{secom\_linear} workloads differ in one component,
\code{cv\_error}, the cross-validation loss evaluated on the threshold grid.
The mean relative difference is $1.40 \times 10^{-16}$ for
\code{secom\_linear\_atlas}, $2.63 \times 10^{-16}$ for
\code{secom\_linear\_2groups} and $3.34 \times 10^{-16}$ for
\code{secom\_linear\_atlas\_genus\_large}; the maximum absolute difference over
the three workloads lies between $1.1 \times 10^{-16}$ and
$8.9 \times 10^{-15}$. The difference arises from the order of summation in the
squared Frobenius norm. The quantities selected from \code{cv\_error},
\code{thresh\_opt} and \code{alpha\_opt}, and the matrices \code{corr\_th} and
\code{corr\_fl} derived from them, are identical.

\code{R CMD check} on the built tarball returns \code{Status: OK}, with 13
passing tests, no failures, no warnings and no skips, and all five vignettes
re-built; the sensitivity analysis table of the rendered \code{ANCOMBC2}
vignette reproduces the counts quoted in its section 3.4.

Table~\ref{tab:scaling} reports configurations outside the harness that isolate
the scaling of individual entry points.

\begin{table}[htbp]
\centering
\small
\caption{Synthetic and auxiliary scaling measurements. Times are single runs;
memory is the master-process peak.}
\label{tab:scaling}
\begin{tabular}{l l r r r r}
\toprule
Function & Configuration & 2.13.2 (s) & 2.15.1 (s) & 2.13.2 (Mb) &
2.15.1 (Mb) \\
\midrule
\code{ancom} & 300 taxa, 100 samples & 93.6 & 0.16 & 164.0 & 62.7 \\
\code{ancom} & 600 taxa, 100 samples & 340.0 & 0.98 & 179.6 & 93.9 \\
\code{ancom} & 1000 taxa, 100 samples & $881^{*}$ & 1.75 & & 155.7 \\
\code{ancombc} & 2000 taxa, 200 samples & 42.1 & 1.39 & 235.4 & 187.9 \\
\code{ancombc2} & 2000 taxa, 200 samples & 75.5 & 4.94 & 327.6 & 247.2 \\
\code{secom\_dist} & 114 taxa, 873 samples, $R = 50$ & 209.4 & 142.2 & 97.0 &
111.4 \\
\bottomrule
\end{tabular}
\end{table}

The synthetic configurations use a negative binomial count table with a
three-level group variable, one continuous covariate and \code{n\_cl = 1}, and
for \code{ancombc2} the global, pairwise and Dunnett tests with
\code{pseudo\_sens = TRUE}. The entry marked $^{*}$ extrapolates the 2.13.2
time from the 300-taxon and 600-taxon times under the fitted power law in the
number of taxa, with exponent 1.86. The \code{secom\_dist} configuration is
atlas1006 at the Genus level with \code{n\_cl = 2}.

\section{Changes by function}

\subsection*{\code{ancom}}

With \code{rand\_formula = NULL}, \code{ancom} builds one design for all
additive log-ratio responses. The model frame, the model matrix, the QR
factorization, the residual degrees of freedom and the inverse cross-product
matrix are computed once by \code{.ancom\_alr\_design} and shared by every pair
of taxa. For one reference taxon, the responses of all remaining taxa form the
columns of a response matrix passed to a single \code{stats::lm.fit} call. In
\code{.ancom\_alr\_stats}, the $F$ statistic of a categorical main variable is
formed from the sequential sums of squares in \code{fit\$effects}, as in
\code{anova.lm}, and the $t$ statistic of a continuous main variable from the
coefficient on the second row of the design and the corresponding diagonal
entry of the inverse cross-product matrix, as in \code{summary.lm}. The
per-pair \code{lm()} path is used when the main variable is not a term of the
model, when the design is rank deficient, when there are no residual degrees of
freedom, or when a response is not finite.

Only the pairs with $j > i$ are fitted; the entry $(j, i)$ is filled by
reflection, with the sign of the $t$ statistic and of the effect estimate
reversed. With \code{rand\_formula} supplied, each pair is fitted by
\code{lme4::lmer}. The \code{foreach} loop iterates over \code{n\_cl}
interleaved blocks of reference taxa and exports the design list and the
log-transformed table; the p-value and effect-size matrices are preallocated.

\subsection*{\code{ancombc} and \code{ancombc2}, fixed effects}

\code{.lm\_fit\_all} groups taxa by their pattern of usable samples, where a
sample is usable for a taxon when its response is finite and its design row is
complete. Each group is solved by one multi-response \code{stats::lm.fit} on
the shared design restricted to the usable samples. A group whose restricted
design is rank deficient is refitted taxon by taxon with \code{stats::lm()}, so
that an absent factor level is dropped in the same way. The pattern keys are
built by a vectorized paste over the columns of the usability indicator.

\code{.sandwich\_vcov} stores the per-sample outer products $x_j x_j^{\top}$,
which do not depend on the taxon, once as the rows of an $n \times p^2$ matrix.
The sum over samples is accumulated for all taxa in a block at a time, in the
sample order used for the per-taxon accumulation, so the floating-point result
is unchanged. The block size bounds the size of the temporaries.

\code{.bias\_em} evaluates \code{dnorm} on the whole coefficient vector in the
E step, and builds the \code{nloptr} option list once per EM run.

The preparation functions compute prevalence by \code{rowSums} in
\code{.data\_core}, build the structural-zero table directly in
\code{.get\_struc\_zero}, draw the $t$ variates of \code{.dunn\_global} for the
whole matrix of degrees of freedom in one call from the same stream, preallocate
the output of \code{.combn\_fun2}, and avoid repeated transposition of the
feature table in \code{.ancombc2\_core}.

\subsection*{\code{ancombc2}, random effects}

\code{.lmer\_struct\_all} groups taxa by their pattern of non-missing responses
and calls \code{lme4::lFormula} once per group. The resulting structure, which
holds the model frame, the fixed-effects model matrix and the random-effects
terms, does not depend on the response values and is reused by every taxon and
every REML iteration in that group.

\code{.lmer\_refit} fits one response from a stored structure by the steps that
\code{lme4::lmer} performs after \code{lFormula}: \code{mkLmerDevfun},
\code{optimizeLmer}, \code{checkConv} and \code{mkMerMod}. The covariance
parameters \code{reTrms\$theta} and the nonzero entries of
\code{reTrms\$Lambdat} are restored to their starting values in a private copy
before each fit, because \code{lme4} modifies them in place during
optimization. Each fit therefore does not depend on the fits that precede it.

\code{lmerTest::lmer} is called only where Satterthwaite degrees of freedom are
required, in the final fit at the converged covariance parameters, and
\code{summary()} is called once per taxon. A taxon whose variance-covariance
matrix is not positive definite is reported as unfitted. The initial pass and
the first REML iteration both evaluate the model at $\theta = 0$, and one pass
serves both. From each fit, the coefficients, fitted values, residuals and
variance-covariance matrix are extracted and the \code{merMod} object is
discarded.

\subsection*{\code{secom\_linear}}

The bodies of the \code{foreach} loops are evaluated in an environment
constructed by \code{.detached\_env}, whose parent is the global environment
and whose contents are the matrices and scalars the body reads, together with
\code{foreach} and \code{\%dorng\%}. The parallel backend therefore loads no
package namespace on the workers at each call. The \code{\%dorng\%} streams are
unchanged.

The co-occurrence matrix is the cross-product of the occurrence indicator
matrix, with its diagonal set to the column sums of that indicator.

\code{.thresh\_loss\_grid} evaluates the cross-validation loss at every point of
the threshold grid in one pass. The entries of the training correlation matrix
are sorted once by absolute value; \code{findInterval} gives the number of
entries set to zero at each threshold; and the loss follows from prefix sums of
$b^2$ and suffix sums of $(a - b)^2$, with the additional suffix sums that soft
thresholding and the weighted penalty require.

The false-positive screen computes the correlation of the estimated
sampling-fraction difference with the columns of the abundance matrix and forms
the mask by \code{outer} on the resulting indicator.

\subsection*{\code{secom\_dist}}

\code{secom\_dist} uses the same detached evaluation environment for its
\code{foreach} loop, the same cross-product form of the co-occurrence matrix and
the same false-positive screen as \code{secom\_linear}. The distance
correlation and its p-value are both taken from the single
\code{energy::dcor.test} call for a pair of taxa.

\section{The \texorpdfstring{\code{conservative}}{conservative} argument}

The sensitivity analysis to pseudo-count addition compares the significance of
a taxon on the complete data with its significance under added pseudo-counts.
The argument \code{conservative} selects the rule, and applies only when
\code{pseudo\_sens = TRUE}. Let $\alpha$ denote the significance level and
$\mathbb{1}\{\cdot\}$ the indicator function. The rules below are applied
column by column to \code{res}, \code{res\_global}, \code{res\_pair},
\code{res\_dunn} and \code{res\_trend}, which carry the same column layout under
both settings.

\paragraph{\code{conservative = TRUE} (default).} The pseudo-count set is
\begin{equation}
\mathcal{P}_1 = \{0.1,\, 0.5,\, 1\}.
\end{equation}
Each value is added to the feature table before the estimation of the sampling
fractions, and the ANCOM-BC2 algorithm is re-run in full. Let $q^{(c)}$ denote
the adjusted p-value of a taxon in the run with pseudo-count $c$, with $q^{(0)}$
the adjusted p-value of the run on the complete data. Then
\begin{align}
\texttt{score} &= \frac{1}{|\mathcal{P}_1| + 1}
  \sum_{c \in \{0\} \cup \mathcal{P}_1}
  \mathbb{1}\{q^{(c)} > \alpha\}, \\
\texttt{passed\_ss} &= \mathbb{1}\{\texttt{score} \in \{0, 1\}\}, \\
\texttt{diff\_robust} &= \texttt{diff} \wedge \texttt{passed\_ss}.
\end{align}

\paragraph{\code{conservative = FALSE}.} The sampling fractions are estimated
once, on the complete data. The pseudo-count set is
\begin{equation}
\mathcal{P}_0 = \{0.01,\, 0.02,\, \ldots,\, 0.50\}, \qquad
|\mathcal{P}_0| = 50 .
\end{equation}
Each value is added to the zero entries of the bias-corrected table, and the
fixed effects are re-estimated by ordinary least squares, one fit per taxon per
pseudo-count, distributed over the \code{\%dorng\%} loop. Let $p^{(c)}$ denote
the unadjusted p-value of a taxon at pseudo-count $c$, and $p^{(0)}$ the
unadjusted p-value of the run on the complete data. Then
\begin{align}
\texttt{score} &= \frac{1}{|\mathcal{P}_0|}
  \sum_{c \in \mathcal{P}_0} \mathbb{1}\{p^{(c)} > \alpha\}, \\
\texttt{passed\_ss} &= \mathbb{1}\{\texttt{score} = 0,\; p^{(0)} \le \alpha\}
  + \mathbb{1}\{\texttt{score} = 1,\; p^{(0)} > \alpha\}, \\
\texttt{diff\_robust} &= \texttt{diff} \wedge \texttt{passed\_ss}.
\end{align}
A taxon with $p^{(0)}$ missing is assigned $p^{(0)} = 1$. This rule is the one
used in Bioconductor release 3.19.

\paragraph{Example counts.} On the QMP simulation of the \code{ANCOMBC2}
vignette (91 taxa), the three \code{res\_pair} comparisons give 7, 7 and 12
taxa with \code{diff\_robust = TRUE} under both settings. The counts of
\code{passed\_ss = TRUE} are 72, 70 and 90 with \code{conservative = FALSE} and
67, 55 and 84 with \code{conservative = TRUE}. On dietswap at the Family level,
the \code{nationalityAFR} column of \code{res} gives 4 taxa with
\code{diff\_robust = TRUE} under \code{conservative = TRUE} and 5 under
\code{conservative = FALSE}; \code{res\_global} gives 3 and 2, and
\code{res\_trend} gives 5 and 3. The direction of the \code{res\_trend}
difference follows from the test used in the refits: the sequential $F$ test of
the ordinary least squares refit is not the constrained trend test of the main
run, and for a repeated-measures design the two disagree on a subset of taxa.
On dietswap, the sensitivity analysis took 239 s with
\code{conservative = TRUE} and 77 s with \code{conservative = FALSE}, measured
on the installation preceding the random-effects changes of
Section~\ref{sec:verification}.

\section{Bug fixes}
\label{sec:bugs}

\paragraph{\code{ancombc} with \code{pseudo = 0}.} The response of the
fixed-effects fit is
\begin{equation}
y_{ij} = \log(o_{ij} + \texttt{pseudo}), \qquad
y_{ij} \leftarrow \texttt{NA} \ \text{ if } y_{ij} \text{ is not finite},
\end{equation}
where $o_{ij}$ is the count of taxon $i$ in sample $j$. With
\code{pseudo = 0} and $o_{ij} = 0$ the first expression gives $-\infty$, and
the second assignment treats that entry as missing, as \code{.ancombc2\_core}
does. The taxon is then fitted on its observed samples. In case 1 of
\code{benchmarks/harness/edge\_cases.R}, a table of 60 taxa and 60 samples with
zero counts in two taxa, three taxa that 2.13.2 returns as \code{NA} are
estimated, and no entry of \code{diff\_abn} changes.

\paragraph{\code{ancombc2} with \code{trend = TRUE} and
\code{global = FALSE}.} The sensitivity analysis of the trend test reads the
sensitivity table of the global test. When \code{trend = TRUE}, the global test
is therefore computed regardless of the value of \code{global}, and its result
is removed from the return value when \code{global = FALSE}. The combination
\code{pseudo\_sens = TRUE}, \code{trend = TRUE} and \code{global = FALSE},
which stopped with \code{'MARGIN' does not match dim(X)}, now returns the trend
results with their \code{passed\_ss} and \code{diff\_robust\_abn} columns and
no \code{res\_global} component.

\section{Remaining costs}

The shares below are from \code{Rprof} profiles of version 2.15.1.

\begin{itemize}
\item In \code{secom\_dist\_atlas}, the compiled permutation loop of
\code{energy::dcor.test}, entered through \code{.C}, accounts for 95.1\% of
self time. The cost is $O(d^2 R)$ permutation replicates for $d$ taxa.
\item In \code{ancombc2\_dietswap}, the random-effects path
\code{.remle\_fit\_all} accounts for 86.2\% of total time, and the \code{lme4}
optimizer called through \code{nloptr} accounts for 39.0\% of total time, of
which the objective function \code{eval\_f} accounts for 30.4\%.
\item In \code{ancombc\_atlas\_genus\_large}, the \code{nloptr} call of the M
step in \code{.bias\_em} accounts for 72.9\% of total time; its objective
function accounts for 54.2\% of total time and 22.9\% of self time, and the
\code{dnorm} evaluations of the E step account for a further 22.9\% of self
time.
\item In the fixed-effects profiles, \code{.sandwich\_vcov} accounts for 4\% to
19\% of self time.
\end{itemize}

\section{Reproducing the measurements}

The scripts take all paths from their command-line arguments. Let
\code{lib\_a} and \code{lib\_b} denote library directories holding the two
installations to compare, and \code{out\_a} and \code{out\_b} the directories
for the results. From the package root:

\begin{verbatim}
Rscript benchmarks/harness/run.R     lib_a out_a all 3
Rscript benchmarks/harness/run.R     lib_b out_b all 3
Rscript benchmarks/harness/compare.R out_a out_b 1e-8
Rscript benchmarks/harness/compare.R out_a out_b 0
\end{verbatim}

\noindent
\code{run.R} takes the library path, the output directory, the scale
(\code{all}, \code{vignette} or \code{large}), the number of repetitions, and
an optional regular expression restricting the workload names. It writes one
\code{.rds} file per workload and \code{timings.csv}. \code{compare.R} takes
the two output directories and the tolerance, and writes
\code{compare\_vs\_<basename(out\_a)>.csv} into the second directory. The edge
cases and the \code{pseudo = 0} reproduction are run as

\begin{verbatim}
Rscript benchmarks/harness/edge_cases.R    lib_b out_b/edge_cases.rds
Rscript benchmarks/harness/repro_pseudo0.R lib_b
\end{verbatim}

\end{document}
